\begin{abstract}
Immersive virtual reality (VR) used in safety-critical settings exposes a perception-layer attack surface: manipulated rendering can induce cybersickness and harm users. We study E-DDQN, a dueling Double Deep Q-Network with an asymmetric reward that penalises missed cybersickness events five times more than false alarms, class-balanced replay and soft target updates, and compare it with a standard DQN and cost-weighted supervised models under an adversarial threat model. Rebuilding the study from the raw public UFFCSData exports (\nSessionsEval{} curated sessions, \nSamplesEval{} one-second samples), we show that previously reported near-perfect accuracy is explained by oversampling before a row-level split: in a reconstruction of that pipeline \LeakDupPct\,\% of positive test rows had exact duplicates in training, and a grouped split reduces the same features to ROC-AUC \LeakGrpRocAuc{}. We prove that on logged data the alarm cannot affect future observations, so the task is a contextual bandit, and that the asymmetric reward induces exactly the 5:1 cost-sensitive Bayes threshold. After correcting the agent (within-session transitions, removal of invalid teacher forcing, target-network statistics), all models were tuned, checkpointed and thresholded on validation groups only within a frozen, nested, participant-conservative protocol. E-DDQN reached a 5:1 expected cost of \EDDQNCostCostMean{} and ROC-AUC \EDDQNCostRocAucMean{}, did not outperform the standard DQN (\DQNCostCostMean, \DQNCostRocAucMean), and was significantly worse than a cost-weighted random forest (\RFCostCostMean, \RFCostRocAucMean; group-bootstrap $p=$\,\PairDOneRFPboot). Only the random forest had a materially lower cost than always alarming (\AlwaysCost), and with \nGroupsEval{} independent groups even its confidence interval includes that cost. Bounded perturbations of frame-rate, orientation and locomotion telemetry changed ROC-AUC by at most \AttAucMaxAbsChange{}, whereas adaptive worst-case evasion lowered recall by \AttEvasionDropMin--\AttEvasionDropMax{} percentage points. The safety objective of near-complete detection with controlled false alarms is not achieved on these data; demonstrating detection of real rendering attacks requires datasets with applied manipulations and measured effects.
\end{abstract}

\begin{IEEEkeywords}
Deep reinforcement learning, double deep Q-network, contextual bandit, cost-sensitive learning, virtual reality, cybersickness, adversarial threat model, data leakage, explainable AI, safety-critical systems.
\end{IEEEkeywords}
